\section{Methodology}
\label{sec:methodology}

\subsection{Motivating Persona Study}

The experiments reported in this paper originate in a finding from the companion study on authoring profile taxonomy~\cite{dellaLibera2025taxonomy}, which evaluated seven bare persona conditions on a controlled puzzle authoring task. Each condition received the instruction ``You are a [role]'' and was given the same task: design an Age of Empires~II feeder puzzle whose final answer is the word TOWER, using only a locked world data file as domain reference. No puzzle design conventions, extraction templates, or hunt design principles were provided. The seven conditions were: null (AX, no role label), puzzle designer (A0), game designer (AG), programmer, doctor, artist (AA), and no job title.

The C8 panel evaluated all seven outputs on the seven-dimension weighted rubric (/45, Clarity ×1 + Solvability ×1 + Elegance ×2 + Reading Reward ×2 + Fun ×1 + Confirmation ×1 + Riven Standard ×1, pass $\geq 33$). Scores ranged from 16.3 (AX-null) to 40.3 (AA-artist), with the game designer condition (AG) second at 32.3 and all other conditions below 33. The artist was the only condition to achieve a passing score and the only condition to produce a mechanism outside the letter-indexing and first-letter acrostic family. All six other conditions, regardless of professional label, extracted TOWER via first-letter or positional indexing. The 15.6-point gap between AA (40.3) and A0 (24.7) is the motivating anomaly this paper explains.

\subsection{Trait Taxonomy}

The six cognitive traits were identified through a two-stage process. In the first stage, we conducted an introspective analysis of the artist persona: given that ``You are an artist'' contains no puzzle domain knowledge, what disposition or orientation could it activate that would drive a model toward experience-first mechanism design rather than template-extraction? The answer we formulated was that artistic training, as documented across design philosophy, creativity research, and practitioner writing, cultivates a small number of domain-agnostic cognitive habits: a primary orientation toward the audience's experience, a drive to produce surprise and recognition rather than mere correctness, a preference for visual and spatial thinking as the first representational step, a discipline of removing excess until the core is exposed, a habit of verification before completion, and a practice of structural coherence before surface elaboration.

In the second stage, each candidate disposition was cross-referenced against the creativity and design theory literature. Experience-first thinking is the operationalization of Norman's user-centered design principle~\cite{norman2013design} and Hassenzahl's experience-design framework~\cite{hassenzahl2010experience}. Surprise instinct connects to Berlyne's theoretical account of aesthetic curiosity and surprise as primary drivers of aesthetic engagement~\cite{berlyne1971aesthetics}. Visual thinking is documented in Arnheim's research on visual thinking as a primary cognitive mode~\cite{arnheim1969visual} and in Shepard and Metzler's foundational mental rotation work~\cite{shepard1971mental}. Elegance drive is the operationalization of Rams~\cite{rams1976design} and Saint-Exup\'ery's reduction principle. Rigor and verification connects to Snyder's construction discipline~\cite{snyder2019craftsmanship}: no puzzle element should be asserted without confirming that it is factually accurate and that the extraction is uniquely determined. Systematic thinking connects to well-documented expertise advantages in structured problem decomposition~\cite{chi1981categorization}.

Each trait was then formulated as a single imperative sentence designed to activate the disposition as a behavioral directive without introducing puzzle domain content. The imperative formulation is deliberate: it specifies a way of working rather than a domain of knowledge, ensuring that any quality gain observed in the trait-isolation conditions is attributable to the cognitive orientation rather than to domain expertise injected via the trait sentence.

The six resulting trait prompts are as follows. T1 (Experience-First): ``You always think first about how the other person will experience what you make.'' T2 (Surprise Instinct): ``You care deeply about surprising people---you want them to pause and think `wait, really?'\,'' T3 (Visual Thinking): ``You think in images and spatial relationships before you think in words.'' T4 (Elegance Drive): ``You remove anything that isn't necessary---you value spare, clean form above all.'' T5 (Rigor/Verification): ``You check your work carefully before calling it done---you never assume it works.'' T6 (Systematic Thinking): ``You approach problems by breaking them into a logical structure before building anything.''

The six traits have a partial mapping onto Boden's three creativity types~\cite{boden2004creative}: T1 (experience-first) and T4 (elegance drive) function as \emph{exploratory} constraints, establishing the conceptual space within which the mechanism is sought; T2 (surprise instinct) and T3 (visual thinking) function as \emph{combinatorial} operators, generating novel pairings within that space; T5 (rigor/verification) functions as \emph{transformational} pruning, eliminating candidates that violate quality gates. T6 (systematic thinking) is structural scaffolding rather than a creativity type per se. This mapping is approximate and the taxonomy makes no claim to exhaustiveness --- other creative orientations may introduce additional traits not represented here.

\subsection{Full Artist Persona}

For Phase~2 (leave-one-out) and as the reference control for all three phases, a full artist persona was constructed by integrating all six trait dispositions into a coherent character description. The integration was designed to avoid a list structure that would make individual traits visible and removable to the model; instead, the traits are woven into a continuous characterization that reads as a natural identity description. The full artist persona used across all leave-one-out conditions is:

\begin{quote}
You are an artist. You always think first about how the other person will experience what you make. You care deeply about surprising people. You think in images and spatial relationships before you think in words. You remove anything that isn't necessary. You check your work carefully before calling it done. You approach problems by building a coherent structure.
\end{quote}

This is the prompt that produced the AA condition (DAMAGED INSCRIPTIONS, 41.7/45 on Phase~2 re-evaluation; published score 40.3/45 from the motivating study, within expected reviewer variance across evaluation sessions). Each leave-one-out condition removes one sentence from this paragraph, replacing the missing trait with no substitute text. The removal is exact: the remaining persona is a strict subset of the full text, preserving all other trait descriptions unchanged.

\subsection{Phase 1: Trait Isolation}

Phase~1 tests each of the six traits in isolation: one trait sentence, one agent, one puzzle authoring task. Each agent receives the bare task (``Design an Age of Empires~II feeder puzzle whose answer is TOWER'') plus the single-sentence trait injection and the locked AoE2 world data file. No persona identity label (``artist,'' ``designer,'' or otherwise) is provided---the trait sentence is the complete persona specification. Conditions T1 through T6 are thus the minimal possible operationalizations of each trait: if any quality gain is observed, it is attributable to the single sentence and nothing else.

Six agents were run in the same session as the AA and A0 controls from the motivating study, providing a shared calibration baseline. All eight puzzles (T1--T6, AA, A0) were read by the C8 panel before any scoring, with an explicit calibration note establishing the mental rank ordering across the full set before any numerical scores were assigned. This panel-reads-all-first protocol is the same calibration procedure used in the companion evaluation papers~\cite{dellaLibera2025panel,dellaLibera2025taxonomy} and ensures that individual puzzle scores are assigned relative to the full distribution rather than evaluated in isolation.

Phase~1 addresses two specific questions. First, which single trait has the greatest independent impact on puzzle quality? Second---and more diagnostic---does any single trait replicate the artist's mechanism novelty? If a single trait (particularly T1 or T3) independently produces a damaged-inscription-type mechanism, the mechanism novelty finding from the motivating study is attributable to a single cognitive orientation rather than to the full artist persona. If no single trait replicates mechanism novelty, the finding is that AA-level performance requires a trait combination.

\subsection{Phase 2: Leave-One-Out}

Phase~2 tests the full artist persona with one trait removed per condition. Six agents run the full paragraph minus one sentence each, producing conditions L-T1 through L-T6. The same task and world data are provided; evaluation uses the same C8 panel and /45 rubric. Because the full persona is the starting point rather than the bare task, Phase~2 measures the \emph{marginal contribution} of each trait to a near-complete creative specification: not what each trait adds from nothing, but what each trait's removal costs.

The leave-one-out design is the standard ablation methodology for identifying load-bearing components in a system where components interact. A trait that produces a large drop when removed is structurally irreplaceable---no other trait or combination of remaining traits compensates for its absence. A trait that produces a small drop when removed is quality-amplifying but substitutable---other traits in the system are sufficient to maintain passing quality even without it.

Phase~2 also answers a diagnostic question that Phase~1 cannot: whether a trait that scores poorly in isolation is in fact load-bearing when embedded in the full artist context. T5 (rigor and verification) scores 27.0/45 in isolation---below passing---but its removal from the full artist persona produces a 12.7-point drop, crossing below threshold. This pattern, shared by T4 (isolation 26.0, removal $-11.7$), indicates that T4 and T5 are not primarily independent quality drivers; they are constraints that prevent the rest of the persona from self-sabotaging. T5 in isolation produces careful but experience-free puzzles; T5 in the full artist context prevents factual errors that would break solvability for the domain-expert panel members who constitute the C8 reviewer population.

\subsection{Phase 3: Cumulative Reconstruction}

Phase~3 starts from the null baseline (AX, no persona, no trait sentences) and adds traits one at a time in tier order as established by Phase~2 results. The reconstruction sequence is: AX $\to$ +T1 $\to$ +T4 $\to$ +T5 $\to$ +T2 $\to$ +T3 $\to$ +T6. Conditions are designated R0 (AX, baseline) through R6 (all six traits, equivalent to the full artist persona). R0 and R6 are in hand from the motivating study and Phase~2 re-evaluation respectively; R1 through R5 are new agents run under the same task and world data conditions.

The build order follows the tier classification implied by Phase~2 leave-one-out results: Tier~1 (irreplaceable, removal $> -20$ points) is added first; Tier~2 (load-bearing, removal $-11$ to $-13$ points) is added second and third; Tier~3 (quality-amplifying but substitutable, removal $-4$ to $-8$ points) is added in order of estimated Phase~2 contribution. This sequencing ensures that the reconstruction trajectory follows the steepest quality gradient---each step adds the trait whose absence was most costly at that tier---and that the first threshold crossing (quality $\geq 33/45$) is identified at the earliest possible reconstruction step.

Phase~3 identifies the minimum sufficient trait set as the smallest R-step set at which quality first crosses the 33/45 pass threshold. All seven conditions (R0 through R6) are evaluated in a single C8 panel session with full calibration, preventing between-session reviewer drift from confounding the reconstruction trajectory. The Phase~3 session result thus provides the cleanest available estimate of how each tier of trait accumulation affects quality, from the floor (AX, 16.3/45 in the motivating study) to the ceiling (AA-equivalent at $\approx$ 41/45).

\subsection{Evaluation Infrastructure}

All experiments use the same evaluation infrastructure as the motivating persona study and the companion reviewer-profile depth paper~\cite{dellaLibera2025panel}: the C8 panel (three domain-expert reviewer profiles, Dan Katz, Thomas Snyder, and Dana Young, instantiated from 80--120 line structured markdown profiles encoding their characteristic evaluative frameworks and review vocabulary), the seven-dimension weighted rubric (/45 total, pass $\geq 33$), and the all-puzzles-before-any-scores calibration protocol. The same AoE2 world data file used in the motivating study provides domain reference for all authoring agents, ensuring that world-knowledge access is held constant across all conditions.

Scores are reported as three-reviewer panel averages, with individual reviewer scores reported in full in the results section for conditions where reviewer spread exceeds two points on any dimension. Because all reported scores represent single evaluation runs at temperature 1.0, scores are point estimates without sampling variance. The three-reviewer design provides a partial proxy for evaluative uncertainty: spread within the panel at any given score dimension signals genuine disagreement that averaging would suppress. The evaluation rubric's double-weighting of Elegance and Reading Reward reflects the finding from~\cite{dellaLibera2025panel} that these two dimensions are the strongest tier discriminators in the hunt quality distribution, and that equal-weight rubrics systematically miscalibrate for creative work that deliberately trades Clarity and Solvability ceiling for Elegance and Reading Reward---precisely the trade-off that distinguishes AA from A0.

A methodological note on evaluation validity: the C8 panel is itself AI-instantiated --- three AI reviewers operating under domain-expert profiles. Using AI judgment to assess AI creative output introduces an epistemological circularity that this study cannot resolve internally. The companion reviewer ablation study~\cite{dellaLibera2025panel} provides the closest available proxy for panel validity (AI panel scores correlated with quality tier separations on 15 real puzzle hunt puzzles), but direct human expert validation of trait-decomposition quality claims remains a required future step.
